\documentclass[12pt]{article}

% Pacotes de Formatação de Texto
\usepackage[portuguese]{babel}
\usepackage[utf8]{inputenc}
\usepackage[a4paper,top=2cm,bottom=2cm,left=2cm,right=2cm]{geometry}
\usepackage{parskip}
\linespread{1.0}

% Pacotes Matemáticos e Gráficos
\usepackage{amsmath}
\usepackage{amssymb}
\usepackage{amsfonts}
\usepackage{cancel}
\usepackage{tikz}

% Outros Pacotes
\usepackage{graphicx}
\usepackage{ascii}
\usepackage{listings}
\usepackage{enumitem}
\usepackage{tcolorbox}
\tcbuselibrary{skins,breakable}
\usepackage{xcolor}
\usepackage{prettyref}

\input{latex_boxes_definitions.tex}

\title{Título}
\author{Nome}
\begin{document}

\title{Exemplos de Caixas de Texto Corrigidos}
\author{Autor}
\date{\today}
\maketitle

\section{Introdução}
Este documento corrigido demonstra o uso das caixas de texto definidas no arquivo separado.

\section{Exemplos}

\begin{CaixaTeo}{Exemplo de Teorema}
    Este é um exemplo de teorema utilizando a caixa personalizada.
\end{CaixaTeo}

\begin{CaixaDef}{Exemplo de Definição}
    Este é um exemplo de definição utilizando a caixa personalizada.
\end{CaixaDef}

\begin{CaixaProp}{Exemplo de Proposição}
    Este é um exemplo de proposição utilizando a caixa personalizada.
\end{CaixaProp}

\begin{CaixaLema}{Exemplo de Lema}
    Este é um exemplo de lema utilizando a caixa personalizada.
\end{CaixaLema}

\begin{CaixaCoro}{Exemplo de Corolário}
    Este é um exemplo de corolário utilizando a caixa personalizada.
\end{CaixaCoro}

\begin{CaixaObs}{Exemplo de Observação}
    Este é um exemplo de observação utilizando a caixa personalizada.
\end{CaixaObs}

\end{document}
